\chapter{Potentials for Quantum Implementation}\label{ch:quantum}

\begin{hon}\Hon
\Proposed\ Nothing in this part is implemented, and nothing in it is
required by any result elsewhere in the book. It is included because the
sheaf Laplacian is exactly the kind of object for which quantum linear
algebra has a story, and because that story has a specific shape worth
recording before it is attempted. Every claim here should be read as a
research programme with an identified failure mode, not as a capability.
\end{hon}

\section{Why this object at all}

The two quantities \textsc{tatiana} computes on every tick are a
Dirichlet energy $\discord(x)=x^\top\Lap x$ and a normalised version of
it. Both are quadratic forms in a sparse, symmetric, positive
semidefinite operator built from local data. This is the standard input
to quantum linear-algebra routines, and the standard caution applies:
the speedups are in the linear algebra, and they are destroyed by the
cost of getting classical data in and scalar answers out.

\section{The Dirac operator and its block-encoding}

\begin{definition}[Sheaf Dirac operator]\label{def:dirac}
Let $\dsheaf=\bigoplus_k\cob^k$ be the total coboundary of the sheaf.
The \emph{sheaf Dirac operator} is
\[
  \Dirac \;=\; \dsheaf+\dsheaf^{\,\dagger},
\]
a self-adjoint operator on $\bigoplus_k \Ck{k}$ satisfying
$\Dirac^2=\Delta$, the total Hodge Laplacian. Its kernel is the space of
harmonic cochains, which by Hodge theory is isomorphic to
$\bigoplus_k \Hk{k}$.
\end{definition}

The programme, stated as a pipeline:
\begin{enumerate}
  \item \textbf{Block-encode $\Dirac$.} Construct a unitary $U$ whose
        top-left block is $\Dirac/\gamma$ for a subnormalisation
        $\gamma\ge\norm{\Dirac}$. For a cellular sheaf with
        \emph{matrix-valued} restriction maps this is the substantive
        step, and the subnormalisation $\gamma$ scales with the
        heterogeneity of those matrices --- which is precisely the
        quantity that makes the sheaf interesting in the first place.
  \item \textbf{Filter onto $\ker\Dirac$.} Apply a polynomial
        approximation to a projector via quantum singular value
        transformation, isolating the harmonic subspace.
  \item \textbf{Extract a scalar.} Estimate $\discord$, or an overlap
        with a prepared state, by amplitude estimation.
\end{enumerate}

\begin{remark}[Where this programme is honest and where it is not]
Step 1 is the only step with genuine novelty for cellular sheaves with
matrix-valued restrictions; steps 2 and 3 are standard. The honest
statement of the advantage is therefore about the \emph{encoding} and
its resource scaling, not about an end-to-end speedup for
\textsc{tatiana}. In the deployed regime the restriction maps are the
identity (Proposition~\ref{prop:b0kernel}), $\Lap=L_\Kc\otimes I_d$, and
the whole object reduces to a graph Laplacian on seven vertices --- for
which classical computation is instantaneous and any quantum treatment
is theatre. The quantum question only becomes non-trivial once
Open Problem~\ref{op:restrictions} is solved and the maps are genuinely
matrix-valued and heterogeneous.
\end{remark}

\section{The obstruction: input and output}

\begin{hon}\Hon
Two costs dominate and both are classical.

\emph{Input.} The stalk data is a set of $384$-dimensional embeddings
produced by a classical encoder. Preparing an amplitude encoding of
these costs, in the absence of a QRAM, time linear in the data --- which
is the same order as computing $\discord$ classically. Any claimed
speedup that assumes free state preparation is assuming away the
problem.

\emph{Output.} $\discord$ is one real number, and amplitude estimation
returns it to precision $\varepsilon$ in $O(1/\varepsilon)$ queries. For
a diagnostic thermometer read once per tick, $\varepsilon$ need not be
small, but neither is the classical cost large. The regime where the
quantum route can win is not ``compute $\rhoscore$ faster''; it is
``compute something about the spectrum of $\Lap$ that is classically
infeasible'', and no such quantity is currently required by the
architecture.
\end{hon}

\section{What would make this real}

\begin{openproblem}[A quantity worth the encoding]\label{op:quantumquantity}
Identify a functional of $\Lap$ that (i) the architecture genuinely
needs, (ii) is classically infeasible at the scale \textsc{tatiana}
reaches, and (iii) is extractable from the block-encoded $\Dirac$ by a
constant number of scalar measurements. Without (i) the programme is
solving an unposed problem; without (ii) it is slower than the classical
route; without (iii) the output cost swamps the encoding.
\end{openproblem}

\begin{openproblem}[Subnormalisation in terms of sheaf data]
Give an explicit bound on the block-encoding subnormalisation $\gamma$
for a cellular sheaf Dirac operator in terms of the restriction-map
norms and the degrees of the complex, and determine how it degrades as
the maps become heterogeneous. This is the resource-scaling statement on
which any advantage claim rests.
\end{openproblem}

\begin{remark}[Relation to the eddy-current programme]
The block-encoding of a cellular sheaf Dirac operator with matrix-valued
restriction maps is the technical core of a separate research programme
on magnetoquasistatic problems, where the sheaf arises from a
discretised PDE rather than from a cognitive architecture and the
subnormalisation scales in material contrast. The mathematics is the
same; the physical problem supplies what
Open Problem~\ref{op:quantumquantity} asks for and \textsc{tatiana},
at present, does not. That asymmetry is the honest reason this part is
short.
\end{remark}
